
\documentclass{article}
\usepackage{colortbl}
\usepackage{makecell}
\usepackage{multirow}
\usepackage{supertabular}

\begin{document}

\newcounter{utterance}

\centering \large Interaction Transcript for game `codenames', experiment `opponent\_easy', episode 0 with reflexion{-}Qwen3.5{-}27B.
\vspace{24pt}

{ \footnotesize  \setcounter{utterance}{1}
\setlength{\tabcolsep}{0pt}
\begin{supertabular}{c@{$\;$}|p{.15\linewidth}@{}p{.15\linewidth}p{.15\linewidth}p{.15\linewidth}p{.15\linewidth}p{.15\linewidth}}
    \# & $\;$A & \multicolumn{4}{c}{Game Master} & $\;\:$B\\
    \hline

    \theutterance \stepcounter{utterance}  
    & & & \multicolumn{2}{p{0.3\linewidth}}{
        \cellcolor[rgb]{0.9,0.9,0.9}{
            \makecell[{{p{\linewidth}}}]{
                \texttt{\tiny{[GM$|$GM]}}
                {'hidden': {'team': ['stadium'], 'innocent': ['england', 'bell', 'degree', 'foot', 'wake', 'bill', 'state'], 'opponent': [], 'assassin': ['mammoth']}, 'revealed': {'team': {'team': [], 'innocent': [], 'opponent': [], 'assassin': []}, 'opponent': {'team': [], 'innocent': [], 'opponent': [], 'assassin': []}}}
            }
        }
    }
    & & \\ \\

    \theutterance \stepcounter{utterance}  
    & & \multicolumn{4}{p{0.6\linewidth}}{
        \cellcolor[rgb]{0.9,0.9,0.9}{
            \makecell[{{p{\linewidth}}}]{
                \texttt{\tiny{[P1$\langle$GM]}}
                \texttt{You have attempted this same task before and did not succeed. These are the plans you wrote after those attempts.} \\
\texttt{Trial \#1: New plan: The previous attempt was aborted because the response included extra text ("You: ") before the required format, violating the instruction to "never include any other text than is required in this form." Additionally, the clue "animal" was too broad and risky, potentially linking to opponent words like "fly," "crane," or "turkey," which would end the turn. Next time, I will output *only* the two required lines starting immediately with "CLUE:" and "TARGETS:", ensuring no prefixes or conversational text. I will also select a more specific clue that strictly connects to my team's words (e.g., "stone" for marble, or "sea" for octopus/fish) while explicitly avoiding any semantic overlap with opponent or kill words.} \\
\\ 
\texttt{Trial \#2: New plan: The previous attempt was aborted because the response included the prefix "You: " before the required format, violating the strict instruction to "never include any other text than is required in this form." Additionally, the clue "stone" was risky as it could semantically link to the opponent's word "cell" (as in a stone cell) or the distractor "foot" (stone foot), potentially causing a turn end. Next time, I will output *only* the two required lines starting immediately with "CLUE:" and "TARGETS:", ensuring no prefixes, conversational text, or markdown blocks. I will also select a highly specific clue that strictly connects to my remaining team words (e.g., "glass" for microscope/marble or "city" for berlin/stadium) while explicitly verifying no semantic overlap with opponent, distractor, or kill words.} \\
\\ 
\texttt{Trial \#3: New plan: The previous attempt failed because I prioritized covering all remaining words over safety, leading to clues that were too broad and triggered opponent guesses. Specifically, "city" linked to the opponent's "square," and "sports" likely linked to the opponent's "cell" (as in a sports cell or team cell) or was simply too generic. Next time, I will adopt a conservative strategy: I will only group words if the clue is extremely specific and has zero semantic overlap with any remaining opponent, distractor, or kill words. If a safe multi{-}word clue cannot be found, I will intentionally target only a single word with a very narrow clue to ensure the turn continues without triggering an opponent guess. I will rigorously cross{-}reference every potential clue against the full list of opponent and distractor words before submitting.} \\
\\ 
\\ 
\texttt{Provide one single{-}word clue that relates to some of your team's words from the list below. You can choose to how many and to which words your clue relates to.} \\
\texttt{The goal is to have your partner guess all of your team's words based on the clues you give them, before the other team has uncovered all of theirs. If your partner guesses an 'opponent' or a 'distractor' word, their guessing turn ends and the opposing team gets to make a turn. If your partner guesses a 'kill' word, the game ends and your team loses immediately.} \\
\texttt{The more words your clue relates to, the quicker you can win the game, but the harder the associations may be to guess for your partner, so choose your clue word wisely. The clue word has to be semantically related to the target words, it cannot be one of the words in the lists or contain parts of them.} \\
\\ 
\texttt{Always give your single{-}word clue and your comma{-}separated list of related target words in the following format and make your answers as short as possible, never include any other text than is required in this form:} \\
\\ 
\texttt{CLUE: <WORD>} \\
\texttt{TARGETS: <TARGETS>} \\
\\ 
\texttt{Your team words are: stadium.} \\
\texttt{Your opponent's team words are: .} \\
\texttt{Distractor words are: england, bell, degree, foot, wake, bill, state.} \\
\texttt{The kill words are: mammoth.} \\
            }
        }
    }
    & \\ \\

    \theutterance \stepcounter{utterance}  
    & \multicolumn{4}{p{0.6\linewidth}}{
        \cellcolor[rgb]{0.9,0.9,0.9}{
            \makecell[{{p{\linewidth}}}]{
                \texttt{\tiny{[P1$\rangle$GM]}}
                \texttt{CLUE: sports} \\
\texttt{TARGETS: stadium} \\
            }
        }
    }
    & & \\ \\

    \theutterance \stepcounter{utterance}  
    & & & \multicolumn{2}{p{0.3\linewidth}}{
        \cellcolor[rgb]{0.9,0.9,0.9}{
            \makecell[{{p{\linewidth}}}]{
                \texttt{\tiny{[GM$|$GM]}}
                \texttt{sports} \\
            }
        }
    }
    & & \\ \\

    \theutterance \stepcounter{utterance}  
    & & & \multicolumn{2}{p{0.3\linewidth}}{
        \cellcolor[rgb]{0.9,0.9,0.9}{
            \makecell[{{p{\linewidth}}}]{
                \texttt{\tiny{[GM$|$GM]}}
                ['stadium']
            }
        }
    }
    & & \\ \\

    \theutterance \stepcounter{utterance}  
    & & & \multicolumn{2}{p{0.3\linewidth}}{
        \cellcolor[rgb]{0.9,0.9,0.9}{
            \makecell[{{p{\linewidth}}}]{
                \texttt{\tiny{[GM$|$GM]}}
                {'word': 'stadium', 'assignment': 'team'}
            }
        }
    }
    & & \\ \\

    \theutterance \stepcounter{utterance}  
    & & & \multicolumn{2}{p{0.3\linewidth}}{
        \cellcolor[rgb]{0.9,0.9,0.9}{
            \makecell[{{p{\linewidth}}}]{
                \texttt{\tiny{[GM$|$GM]}}
                \texttt{opponent has won} \\
            }
        }
    }
    & & \\ \\

    \theutterance \stepcounter{utterance}  
    & \multicolumn{4}{p{0.6\linewidth}}{
        \cellcolor[rgb]{0.95,0.95,0.95}{
            \makecell[{{p{\linewidth}}}]{
                \texttt{\tiny{[P2$\rightarrow$P1]}}
                \texttt{Post{-}episode Reflexion (not sent to GM):} \\
\texttt{New plan: The previous attempt failed because the clue "sports" was too broad, likely triggering a guess for the distractor "foot" (as in "football") or "bell" (as in a sports bell), causing the turn to end prematurely. Although the format was correct, the semantic safety was insufficient. Next time, I will select a clue that is hyper{-}specific to the target "stadium" while having zero semantic overlap with the distractors (england, bell, degree, foot, wake, bill, state) and the kill word (mammoth). I will choose a concrete noun like "seats" or "field" which strongly implies "stadium" but is unlikely to be associated with the specific distractors provided, ensuring the partner guesses the target without triggering a turn{-}ending error.} \\
            }
        }
    }
    & & \\ \\

\end{supertabular}
}

\end{document}
